\documentclass[10pt,letterpaper]{article}
\usepackage[margin=0.8in]{geometry}
\usepackage[T1]{fontenc}
\usepackage{lmodern,microtype,enumitem}
\usepackage[hidelinks]{hyperref}
\setlength{\parindent}{0pt}
\setlength{\parskip}{5pt}
\setlist[itemize]{leftmargin=*,itemsep=2pt,topsep=3pt}
\begin{document}
\setcounter{page}{4}
\begin{center}
{\Large\bfseries Draft-paper submission notes}\\[3pt]
Reviews in Context or in Weights?\quad 28 September 2026
\end{center}

\textbf{Model paper.}
Amanda Bertsch et al., \emph{In-Context Learning with Long-Context Models: An In-Depth Exploration},
\href{https://arxiv.org/html/2405.00200v2}{arXiv:2405.00200v2} (2025).
Its argument compares contextual demonstrations and fine-tuning across data budgets.
We retain the empirical comparison, experimental setup, diagnostic analysis, and
bounded conclusion. Our evidence concerns reviewer corrections and incomplete
pilots, so the revised comparison is prospective.

\textbf{Adapted outline and visual roles.}
\begin{itemize}
\item \textbf{Abstract and introduction:} state the storage question, conditional
thesis, and two pilot obstacles. Establish that validity and outcome coverage are
necessary to interpret a threshold. Avoid implying that the proposed comparison
has been executed.
\item \textbf{Related work:} place demonstration retrieval, learned playbooks,
low-rank adaptation, self-distillation, and reviewer reliability together.
Move related work earlier than the model paper to establish the playbook comparator.
Keep this compact to preserve room for methods and evidence.
\item \textbf{Completed setup:} describe the agent, editor, public/private checks,
training recipe, held-out split, repeats, and score denominators. This performs
the role of the model paper's experimental setup, with two task-specific metrics.
\item \textbf{Audited results:} separate Lean's failed outcome calibration from
writing's validity-adjusted learning curve and matched-archive controls.
\textbf{Table 1 is the central visual}: it shows every measured writing budget,
usable target count, and score. A table replaces a smooth curve to expose the
change in budget definition and avoid implying an interpolated crossover.
The text gives endpoint uncertainty and validity alongside conditional style.
\item \textbf{Revised experiment:} add the shared-pool arms, reviewer and task
gates, primary endpoint, and prospective inference. This is a necessary addition
because the pilots do not resolve the main thesis. No synthetic result figure
stands in for the unexecuted experiment.
\item \textbf{Discussion and conclusion:} condense mechanism claims into observable
limitations and the remaining empirical question. Omit attention-pattern and
input-ordering analyses from the model paper because these pilots did not test
them. Omit its dataset survey and multi-model figures; our single table and
task-specific audit serve this draft's narrower evidence.
\item \textbf{References:} seven primary research sources. Detailed source
locators, computations, and every claim's justification remain in the separate
\texttt{CLAIMS.md} audit rather than inflating the paper.
\end{itemize}

\textbf{Revisions and status.}
The starting artifact was the 799-word v3 proposal at commit
\texttt{db853b1}, not a prior full class submission.
This draft adds a literature synthesis, completed-pilot methods and results,
limitations, references, and a prospective analysis section.
It separates live and shared-archive comparisons, retains the withdrawal of the Lean
crossover claim, distinguishes $U S$ from valid-only style, and replaces the
``upper bound'' label for fresh-base SFT with ``joint-training reference.''
No instructor or peer feedback was supplied, so none is represented as addressed.

\textbf{Length and audit.}
The paper is intentionally half length under the user's instruction; its measured
word and page counts are in \texttt{evidence/build-check.json}.
This combined PDF contains the draft followed by these notes, as requested by the
\href{https://web.stanford.edu/class/cs197/assignments/project.html\#draft-paper}{CS197 assignment}.
The separate claim ledger classifies observations, literature, design choices,
inferences, and scope limits. The revised run, human checks, and peer-review
exercise remain future work; no new training or inference was performed for this draft.
\end{document}
